\documentclass[10pt]{article}
\usepackage{vmargin}
\setpapersize{USletter}
\usepackage{ragged2e}
\usepackage{amsmath,amssymb}
\usepackage{graphicx}	% Including figure files
\usepackage[caption=false]{subfig}
\graphicspath{{ResponeFigures/}}
\usepackage{hyperref}
\usepackage[shortlabels]{enumitem}

\hypersetup{
	colorlinks   = true, %Colours links instead of ugly boxes
	urlcolor     = teal, %Colour for external hyperlinks
	linkcolor    = teal, %Colour of internal links
	citecolor   = red %Colour of citations
}
\usepackage{xcolor}
\author{T.Kimpson, S. Suvorova, H. Middleton, C. Liu, A. Melaots, R. Evans, W. Moran}
\title{Response to referee comments: \\ PRD Submission DM13138,}
\begin{document}
	\maketitle
\noindent Author response to the referee report for the paper DM13138,
entitled \textit{`Adaptive cancellation of mains power interference in continuous gravitational wave searches with a hidden Markov model'}. \newline 

\noindent We thank the referee for their constructive review.
We agree with all the comments and have made the recommended changes as described below. The referee comments are reproduced in italics for convenience. We separate the response into ``main" comments (indexed alphabetically) and ``additional" comments (indexed numerically), following the convention of the referee.  \newline 



\newpage
\noindent \Large \textbf{Referee comments} \normalsize

\vspace{5mm}

\noindent \textbf{Main comments}

\begin{enumerate}[A.]
	\item \textit{In section II, it is not sufficiently clear if the model is  
		intended as a realistic representation of the North American power grid, or just  
		a phenomenological model that is "good enough" for the purpose of subtraction.  
		This should be spelled out better. An authoritative reference on the actual  
		grid's design and/or statistical properties, if the authors can find one, would  
		also be highly valuable for the typical target audience of this article  
		(astrophysicists / data analysts) who are unlikely to have such information at  
		hand.}
	
		\textbf{Author response:} 
	
	Yes, we agree with the referee that it is important to specify clearly that the model is intended as a phenomenological model that seeks to reproduce the main qualitative features of the mains power (i.e. constant central frequency with small amplitude stochastic wandering, see third paragraph in Section IIB.). We have followed the referee's suggestion and updated the discussion in Section II to make this point explicitly.
	
	 \textcolor{red}{TK: I cannot find any references for the empirical behaviour of the grid. Perhaps the EE-side might know.}
	
	
	
		\item \textit{In the end, it is not clear enough how big the positive impact of  
			the method really is. Especially, and regardless of the answer to point A, it  
			seems necessary to present at least some level of real-data test of the proposed  
			method, beyond the presented simulated noise case. Since LIGO data is readily  
			available via the GWOSC site, this should be straightforward to perform. A  
			clearer summary of the results in abstract and conclusion would then be  
			appreciated.}
	
	\textbf{Author response:} 
	


Yes we agree that a clearer summary of the main results quantifying the impact of the method is important to include. We have updated the abstract and the conclusion accordingly. Specifically
\begin{itemize}
	\item For the representative system (Table 1) ANC suppresses the 60-Hz interference by approximately 20 dB.
	\item Without ANC, the HMM tracker is unable to track the wandering CW signal. After ANC, the HMM tracker is able to track the wandering CW signal. The time-averaged RMS error in the HMM tracking after ANC is 0.038Hz for low $\sigma_f$ and $0.47 Hz$ for high $\sigma_f$.
	\item The detection probability at a 5\% false alarm rate, $p_{\rm d}(0.05) \approx 0.5 $--$0.6$ across a range of mains power interference parameters and filter parameters.
\end{itemize}


\textcolor{red}{TK: Tests on real-data to be discussed in person. A summary of some thoughts below.}

We currently generate some synthetic strain-channel data $x(t) = h(t) + c(t) + n(t)$, and some synthetic reference-channel data $r(t)$, according to some analytical expressions which are specified in the manuscript.

Which of these components would be appropriate to replace with real GWOSC data?

We don't have GWOSC $h(t)$ - no CWs have been detected! Let's assume that in some strain data $x(t)$ that we obtain from GWOSC that there is no $h(t)$ signal in there. The procedure for testing the method on real data might then look like:

\begin{itemize}
	\item Take some strain-channel data $x(t)$ from GWOSC
	\item Simulate some synthetic $h(t)$ according to the manuscript equations. 
	\item Inject the $h(t)$ into $x(t)$ to obtain some data $y(t) = h(t) + x(t)$ which gets passed to the ANC scheme.
	\item Take $N \leq 9$ PEM channels from GWOSC. Pass these to the ANC scheme as $r_{\rm empirical}(t)$.
	\item Run ANC scheme - can we recover our injected $h(t)$? 
\end{itemize}


	
	
	
\end{enumerate}



\noindent \textbf{Additional comments}




\begin{enumerate}
	\item \textit{The abstract currently points out that the validation test is on "synthetic  
		signals", which is perfectly fine, but what is more important is the nature of  
		the noise data.  }
	
	\textcolor{red}{TK: this is related to point B, re real data. To be updated once point B has been answered}
	
	\item \textit{A brief summary of the results would be valuable at the end of the abstract  
		(as per point B above). }
	
	\textbf{Author response:}  \newline 
	Yes, we agree that the main results should be better summarised in the abstract. We have updated the abstract accordingly.
	
	\item \textit{A few comments on references the introduction section:  
		The first 3 are very old, describing glitches in initial detectors - consider  
		updating to 2G references like the Davis+ papers already cited later.  
		The standard Advanced Virgo detector reference arxiv:1408.3978 seems to be  
		missing.  
		Please check if [21] (Lee et al) was indeed the intended citation - this seems  
		to be a deep learning waveform model paper, not about artifact identification.  
		At the end of the second paragraph, another important class of line vetoes not  
		cited yet are those based on multi-detector consistency (e.g.  
		doi:10.1088/0031-8949/90/12/125001 and doi:10.1103/PhysRevD.89.064023)}
	
	
		\textbf{Author response:}  \newline 
	We thank the referee for bringing to our attention references which we had previously overlooked, and the erroneous [21] Lee et al. citation. We have now included all the suggested references, removed reference [21] , and included discussion on the class of vetos based on multi-detector consistency
	
	\item \textit{In the paragraph below, the statement that candidates are vetoed blindly  
		without checking the strength of the noise contribution is not true for all  
		standard CW pipelines. More nuanced discussions can be found in some recent LVK  
		results papers.}
	
			\textbf{Author response:}  \newline 
			Yes we agree with the referee that the previous statement was too broad. We have now updated the paragraph to caveat that some pipelines do have ``line robust" algorithms  which combine statistics from different detectors to suppress single detector artifacts (e.g. Keitel et al. 2015, Keitel 2016a,b).
%			 https://arxiv.org/abs/2201.00697
	
	\item \textit{Please try using multiple linestyles for better readability of figures  
		1,6,8,9. In particular, in Fig. 6 it is difficult to discern that the blue x(f)  
		curve actually shows the signal peak, as it is mostly hidden behind the green  
		h(f).  }
	
	\textbf{Author response:}  \newline 
	We have updated all linestyles in the figures to improve readability, as suggested.
	
	
	\item \textit{At the end of II.A, there should be an astrophysical reference for the Crab's  
		frequency, not just a software paper. (Or both together, since citing software  
		used is indeed commendable.) Since the Crab's frequency evolves relatively  
		quickly, a reference year would also be appropriate when quoting it with decimal  
		digits. }
	
	
		\textbf{Author response:}  \newline 
		We have updated the reference to the Crab frequency to include both an astrophysical reference and a software reference. We now also quote the Crab frequency for a given reference year, as suggested.
	
	\item \textit{Fig. 1: missing definition of the abbreviation "ASD"  }
	
		\textbf{Author response:}  \newline 
The abbreviation ASD is now defined in the caption of Figure 1.
	
	
	\item \textit{After Eq. (1) the strict separation of frequency and amplitude modulation is  
		misleading, as both parts of the Earth's motion contribute to both modulations.}
	
		\textbf{Author response:}  \newline 
The text has been updated to remove the misleading separation between the two modulations.


	
	\item \textit{Before Eq. (4), when talking about the coupling "in various complicated  
		ways", would this not naturally be expected to cause stochastic overlaps of  
		components with different time delays tau? Could the authors comment on whether  
		this is a likely explanation for the observed losses in coherence?  }
	
	\item \textit{footnote 3: gitlab file URLs can be quite ephemeral. Is there no better  
		(more stable) way to reference this list, such as a GWOSC release with DOI, a  
		LIGO DCC entry, etc?}
	
	\item \textit{The statements at the end of II.C and in the Fig. 4 caption contradict each  
		other: are the higher-frequency peaks harmonics of the power mains, or "of  
		different provenance"?}
	
	\item \textit{ In the coherence discussion, it remains unclear what threshold is considered  
		to be still sufficient, and below which the channel would no longer be  
		considered a reliable witness. }
	
	\item \textit{In Eq. (12), it is not made explicit what argument the "arg min" is taken  
		over. Would it be appropriate to write $e_k(w)$ with an explicit argument?}
	
	
	\item \textit{In IV.A, can the authors make an argument why neglecting the Earth's  
		movement is not changing qualitatively the signal-vs-line behaviour? Since CW  
		signals are expected to modulate very slowly intrinsically, typically the shape  
		and extent of their detector-frame frequency spectrum should be dominated by  
		these terms.  }
	
	\item \textit{Also in the same discussion, there is a reference to the F-statistic. This  
		is taking a maximum likelihood over amplitude parameters, not over those related  
		to the Earth's movement. Indeed F-statistic searches typically also include a  
		"demodulation", which is probably what the authors had in mind here. But the way  
		it is currently phrased will confuse a non-expert reader into thinking it  
		maximizes over frequency evolution parameters.  }
	
	
	\item \textit{In IV.B, the definition of signal-to-noise ratio does not match the standard  
		in the CW literature, which usually includes the observation time. The duration  
		is also important for the frequency resolution. Is it the 800s visible in Fig.  
		7? Please clarify more explicitly. }
	
	
	\item \textit{The case (ii) "not plotted for brevity" would seem relevant to show - since  
		now fgw lies within the affected frequency range, it sounds like it could be  
		qualitatively different, and more impressive to demonstrate that the method can  
		still deal with it?  }
	
	
	\item \textit{At the end of IV.B, have the authors verified that the Earth Doppler effect  
		is negligible compared to these frequency errors?  }
	
	\item \textit{In figures 8/9, it would be interesting to also compare with the performance  
		without ANC - is that similar to, or even worse than the "random guess"?  }
	
	
	\item \textit{Why is the "zero interference ROC curve" shown only in Fig. 9, not 8? Is it  
		understood why the best performing ANC results are still quite far away from  
		this, if before the subtraction looked visually very impressive? }
	
	
	\item \textit{Near the end of V.B, latency issues are discussed. Are these really relevant  
		for CW searches, which are usually run on archival data?}
\end{enumerate}




\end{document}